\documentclass[11pt]{article}
\usepackage[review]{acl}
\usepackage{times}
\usepackage{latexsym}
\usepackage{amsmath,amssymb,amsfonts}
\usepackage{graphicx}
\usepackage{booktabs}
\usepackage{multirow}
\usepackage{array}
\usepackage{microtype}
\usepackage{inconsolata}
\usepackage{hyperref}
\usepackage{placeins}
\providecommand{\citename}[1]{#1}
\usepackage{ifthen}
\usepackage{todonotes}

\newboolean{reviewmode}
\setboolean{reviewmode}{true}   % flip to false to hide all notes

\title{A Benchmark for Sinhala-Script Language Identification: \\ Disambiguating Sinhala, Pali, and Sanskrit in Closed and Open-World Contexts}

\author{Anonymous ACL submission}

\begin{document}
\maketitle

\begin{abstract}
Multilingual language identification (LangID) models typically handle script sharing well for international writing systems like Latin or Cyrillic. However, for scripts used primarily within a single country, systems often treat the script as an exclusive proxy for the national language. For the Sinhala script (\texttt{U+0D80}--\texttt{U+0DFF}), this assumption overlooks canonical Pali and classical Sanskrit, which have been written in this script for over two millennia. Consequently, leading foundation models---including fastText, GlotLID, and NLLB---show \textbf{0.0\% zero-shot recall} on Sinhala-script Pali and Sanskrit, misclassifying all passages as modern Sinhala. In this paper, we present the first dedicated LangID benchmark and empirical study for these languages in the Sinhala script across closed-world and open-world settings. On a document-stratified corpus of 74,318 sentences, seven baseline models achieve over 0.98 Macro-F1 on full sentences, but single-word tests reveal that character $n$-gram models substantially outperform neural sequence models ($0.858$ vs.\ $0.610$) by capturing inflectional affixes. In open-world benchmarks across FLORES+, CommonLID, and WiLI-2018, naive fine-tuning causes severe catastrophic forgetting in XLM-RoBERTa (Macro-F1 drops to $0.117$), whereas rehearsal with multilingual data (Aya) restores performance ($0.874$ for XLM-R; $>0.994$ for GlotLID and NLLB).
\end{abstract}

\section{Introduction}
Language Identification (LangID) is a fundamental first step in multilingual text processing. It filters web crawls, routes incoming text to translation systems, and categorizes digital library collections \cite{jauhiainen2019automatic,lui-baldwin-2012-langid,caswell-etal-2020-language}. While existing models successfully handle script-sharing across languages that use Latin or Cyrillic, they often struggle when a script is primarily associated with a single country. In such cases, models tend to treat the script as an exclusive proxy for that country's dominant language \cite{kargaran-etal-2023-glotlid,costa2022no}.

This assumption fails for the Sinhala script (\texttt{U+0D80}--\texttt{U+0DFF}). While the script is primarily used for \textbf{Sinhala}, an Indo-Aryan language spoken by approximately 17 million people in Sri Lanka, it has also been used for over two millennia to write canonical \textbf{Pali}---the liturgical language of Theravada Buddhism---and classical \textbf{Sanskrit} \cite{Walshe_1983}. In fact, the Theravada Buddhist canon (\textit{Tipitaka}) was first committed to writing in the Sinhala script during the first century BCE at the Aluvihara temple \cite{adiga2006ancient}. Consequently, thousands of historical palm-leaf manuscripts (\textit{ola}) and modern printed editions are written entirely in the Sinhala script.

Because current LangID models---such as fastText LID-176 \cite{joulin-etal-2017-bag}, GlotLID v3 \cite{kargaran-etal-2023-glotlid}, NLLB LID-218 \cite{costa2022no}, OpenLID-v2 \cite{burchell-etal-2023-open}, ConLID \cite{foroutan-etal-2026-conlid}, and XLM-RoBERTa \cite{conneau-etal-2020-unsupervised}---map the Sinhala Unicode block exclusively to modern Sinhala, they show \textbf{0.0\% recall} on Sinhala-script Pali and Sanskrit (Table \ref{tab:open_world_full}). Every Pali chant, Buddhist verse, and Sanskrit inscription written in Sinhala script is misclassified as modern Sinhala, polluting multilingual corpora and digital archives.

In this paper, we present the first dedicated LangID benchmark and empirical study for Sinhala, Pali, and Sanskrit in the Sinhala script across closed-world and open-world settings. Specifically, we address the following questions:
\begin{enumerate}
    \item How well do standard machine learning and neural architectures distinguish these three languages when input length decreases from full sentences to single words?
    \item Why cannot lightweight classical models simply be scaled to classify hundreds of global languages?
    \item How can pretrained foundation models be adapted to identify these languages without suffering catastrophic forgetting on other world languages?
\end{enumerate}

To answer these questions, we construct a document-stratified corpus of 74,318 sentences, evaluate seven baseline models across four input lengths, analyse the scalability limits of classical classifiers, and benchmark six foundation models on three 11-language hybrid test suites (FLORES+, CommonLID, and WiLI-2018).

\section{Dataset and Experimental Setup}
\label{sec:dataset}

\subsection{Target Corpus Construction}
A common issue when evaluating text classifiers on historical or religious literature is data leakage from random sentence splits. In canonical Buddhist texts and historical chronicles, formulaic greetings, recurring epithets, and standardized verses appear across multiple sentences. If sentences from the same text are shuffled randomly into train and test sets, models can achieve high scores simply by memorizing specific phrases rather than learning genuine grammatical distinctions. To prevent this, we split our data strictly by document (81.1\% Train, 9.4\% Val, 9.5\% Test; Table \ref{tab:target_corpus}):
\begin{itemize}
    \item \textbf{Sinhala (\texttt{Sinh-Sinh}, 32,263 sentences):} Drawn from the SiPaKosa repository (historical and biographical texts) alongside the Sinhala side of a sentence-aligned Pali--Sinhala parallel corpus.
    \item \textbf{Pali (\texttt{Pali-Sinh}, 29,317 sentences):} Sourced from canonical Theravada \textit{Tipitaka} treatises, \textit{Atthakatha} commentaries, and the SiDiaC-v2.0 corpus \cite{jayatilleke-de-silva-2025-sidiac}.
    \item \textbf{Sanskrit (\texttt{San-Sinh}, 12,738 sentences):} Consisting of the historical 19th-century \textit{SansinNT} (7,953 sentences); classical poetry from the Digital Corpus of Sanskrit (DCS; 4,684 sentences) \cite{hellwig2010dcs}, transliterated into Sinhala script via \textit{Aksharamukha} phonetic mapping \cite{rajan2023aksharamukha}; and 101 protective stanzas (\textit{kavaca}) from SiDiaC-v2.0.
\end{itemize}

\begin{table}[t]
\centering
\small
\setlength{\tabcolsep}{3.5pt}
\caption{Target corpus statistics ($N=74,318$). Partitioned with document-level grouping to prevent data leakage.}
\label{tab:target_corpus}
\resizebox{\columnwidth}{!}{%
\begin{tabular}{lrrrr}
\toprule
\textbf{Language Class} & \textbf{Train} & \textbf{Val} & \textbf{Test} & \textbf{Total} \\
\midrule
Sinhala (\texttt{Sinh-Sinh}) & 26,675 & 2,895 & 2,693 & 32,263 \\
Pali (\texttt{Pali-Sinh})    & 23,490 & 2,800 & 3,027 & 29,317 \\
Sanskrit (\texttt{San-Sinh}) & 10,120 & 1,291 & 1,327 & 12,738 \\
\midrule
\textbf{Total}               & \textbf{60,285} & \textbf{6,986} & \textbf{7,047} & \textbf{74,318} \\
\bottomrule
\end{tabular}%
}
\end{table}

\subsection{Multilingual Evaluation Suites}
In practical settings, Sinhala-script text appears alongside regional neighbors and international languages. We adopted the standardized \textbf{BCP-47 taxonomy} \cite{phillips2009tags} (\texttt{[Language]-[Script]}):
\begin{itemize}
    \item \textbf{Target Triad:} \texttt{Sinh-Sinh}, \texttt{Pali-Sinh}, \texttt{San-Sinh}.
    \item \textbf{Regional Neighbors:} \texttt{San-Deva} (Devanagari Sanskrit), \texttt{Hin-Deva} (Hindi), \texttt{Ben-Beng} (Bengali), and \texttt{Tam-Taml} (Tamil).
    \item \textbf{Global Benchmarks:} \texttt{Eng-Latn} (English), \texttt{Fre-Latn} (French), \texttt{Ger-Latn} (German), and \texttt{Ara-Arab} (Arabic).
\end{itemize}
We constructed three hybrid evaluation suites by mixing our 7,047 held-out target test sentences into the official test splits of \textbf{FLORES+} \cite{goyal-etal-2022-flores}, \textbf{CommonLID} \cite{suarez-etal-2026-commonlid}, and \textbf{WiLI-2018} \cite{thoma2018wili}.

\section{Closed-World Evaluation and Scalability}
\label{sec:closed_world}

\subsection{Model Comparison across Input Lengths}
We began by training seven standard architectures from scratch on our target training set (Table \ref{tab:target_corpus}): Multinomial Naive Bayes (MNB), Linear SVM, Logistic Regression (LogReg), XGBoost, Char-CNN, Char-BiGRU, and fastText (details in Appendix \ref{sec:appendix_architectures}).

\begin{table}[t]
\centering
\small
\setlength{\tabcolsep}{3pt}
\caption{Macro-F1 of baseline architectures trained from scratch across varying input lengths on held-out target test data ($N=7,047$).}
\label{tab:stress_testing}
\resizebox{\columnwidth}{!}{%
\begin{tabular}{lcccc}
\toprule
\textbf{Architecture} & \textbf{Full Sent.} & \textbf{5-Word} & \textbf{3-Word} & \textbf{1-Word} \\
\midrule
Multinomial NB        & 0.9947 & 0.9782 & 0.9520 & \textbf{0.8584} \\
fastText (Scratch)    & 0.9958 & 0.9768 & 0.9416 & 0.8027 \\
Char-CNN              & 0.9950 & 0.9749 & 0.9388 & 0.7654 \\
Char-BiGRU            & 0.9948 & 0.9710 & 0.9255 & 0.6916 \\
Linear SVM            & \textbf{0.9960} & \textbf{0.9801} & \textbf{0.9532} & 0.6099 \\
Logistic Regression   & 0.9955 & 0.9780 & 0.9482 & 0.6054 \\
XGBoost               & 0.9842 & 0.9554 & 0.9102 & 0.5891 \\
\bottomrule
\end{tabular}%
}
\end{table}

\paragraph{Sentence-Level Classification.}
On complete sentences, every model scores above $0.984$ Macro-F1 (Table \ref{tab:stress_testing}), led by Linear SVM at $0.9960$. This indicates that when full sentence context is available, separating these three languages is straightforward for standard classifiers.

\paragraph{Performance on Short Fragments.}
In practical workflows---such as web queries, catalog lookups, or fragmented OCR lines from damaged palm leaves---input text often consists of only a few words. To test how resilient these models are under reduced context, we truncated each test sentence down to its first 5 words, 3 words, and single words.

As shown in Table \ref{tab:stress_testing}, performance diverges noticeably on single words. Linear SVM drops from $0.9960$ to $0.6099$ (a 38.6-point drop), Char-BiGRU falls to $0.6916$, and XGBoost dips to $0.5891$. In contrast, \textbf{Multinomial Naive Bayes maintains an F1 of 0.8584}, and \textbf{fastText achieves 0.8027}.

This difference stems from how models process morphological markers. Deep sequence models and margin classifiers rely on broad contextual patterns. In isolated words, this context is missing. However, Sinhala, Pali, and Sanskrit are highly inflected languages with distinctive affixes:
\begin{itemize}
    \item \textbf{Pali} preserves distinctive nominal case endings (genitive \textit{-assa}, instrumental plural \textit{-\={e}hi}) and verbal suffixes (\textit{-ti}, \textit{-ntu}).
    \item \textbf{Sanskrit} features heavy consonant clusters, visarga endings, and diagnostic dual/ablative markers (\textit{-\={a}t}, \textit{-bhy\={a}m}).
    \item \textbf{Sinhala} exhibits modern postpositions and colloquial verbal suffixes (\textit{-v\={e}}, \textit{-\={e}ya}, \textit{-da}).
\end{itemize}
Because Multinomial Naive Bayes and fastText directly model character $n$-gram frequencies, they can identify these sub-word inflectional markers even in isolated words without surrounding grammar.

\subsection{Scalability Limitations of Classical Classifiers}
Given that simple models like Naive Bayes and Linear SVM achieve over $0.99$ Macro-F1 on full sentences, an obvious question is: \textit{Why not simply train an SVM or Naive Bayes classifier on all 176+ world languages?} We identify four practical and mathematical reasons why this is not feasible:

\paragraph{1. Feature Space Growth $\mathcal{O}(|V|)$.}
Within a single script, character 1-to-5-grams generate around $45,000$ unique features ($|V_{\text{target}}| \approx 4.5 \times 10^4$). However, across $K \ge 176$ languages spanning more than 30 distinct Unicode writing systems (Latin, Cyrillic, Arabic, Devanagari, Hanzi, etc.), the vocabulary of character $n$-grams exceeds $|V_{\text{global}}| > 2.5 \times 10^7$ dimensions. For kernel SVMs, computing kernel matrices scales as $\mathcal{O}(N^2 \cdot |V|)$, which exhausts memory when training on millions of documents.

\paragraph{2. Linear Classification Overhead $\mathcal{O}(K \cdot D)$.}
Classical models rely on flat Softmax or One-vs-Rest setups. Calculating probabilities over $K \ge 176$ classes requires $\mathcal{O}(K \cdot D)$ operations per sample. In contrast, modern foundation models like fastText avoid this bottleneck by using \textbf{Hierarchical Softmax} organized as a balanced Huffman tree \cite{joulin-etal-2017-bag}, cutting per-token inference down to $\mathcal{O}(\log_2 K)$, while metric-based models like ConLID map text into compact continuous representations \cite{foroutan-etal-2026-conlid}.

\paragraph{3. Lack of Subword Representation Sharing.}
Classical $n$-gram representations treat every subword feature as an independent, orthogonal dimension. They have no way to recognize that words with shared Sanskrit roots in Hindi, Bengali, Pali, and Sinhala carry related morphological meaning. Dense neural and subword models embed tokens into continuous vector spaces ($\mathbb{R}^d$), naturally transferring morphological similarities across related language families.

\paragraph{4. Corpus Collection and Balancing.}
Training a competitive classical model from scratch across 200 languages requires gathering, cleaning, and balancing hundreds of millions of web pages. Pretrained foundation models (fastText LID-176, GlotLID, NLLB) have already absorbed massive web crawls. Adapting these models to resolve local ambiguities is far more practical than re-crawling the web.

\section{Open-World Multilingual Evaluation}
\label{sec:open_world}

\subsection{Zero-Shot Performance and Fine-Tuning}
We evaluated six prominent foundation models on our 11-language hybrid test suites (Table \ref{tab:open_world_full}): \textbf{fastText LID-176}, \textbf{OpenLID-v2}, \textbf{NLLB LID-218}, \textbf{GlotLID v3}, \textbf{ConLID}, and \textbf{XLM-RoBERTa Base}.

\begin{table*}[t]
\centering
\small
\setlength{\tabcolsep}{3.5pt}
\caption{Macro-F1 across 11 languages on FLORES+, CommonLID, and WiLI-2018 ($N=7,047$ target + background test sets) under Zero-Shot, Target-Only Fine-Tuned, and Continual Multilingual Rehearsal (Aya) conditions.}
\label{tab:open_world_full}
\resizebox{\textwidth}{!}{%
\begin{tabular}{l ccc ccc ccc}
\toprule
\multirow{2}{*}{\textbf{Foundation Model}} & \multicolumn{3}{c}{\textbf{Zero-Shot Pretrained}} & \multicolumn{3}{c}{\textbf{Target-Only Fine-Tuned}} & \multicolumn{3}{c}{\textbf{Continual Rehearsal (Aya)}} \\
\cmidrule(lr){2-4} \cmidrule(lr){5-7} \cmidrule(lr){8-10}
& \textbf{FLORES+} & \textbf{CommonLID} & \textbf{WiLI-18} & \textbf{FLORES+} & \textbf{CommonLID} & \textbf{WiLI-18} & \textbf{FLORES+} & \textbf{CommonLID} & \textbf{WiLI-18} \\
\midrule
XLM-RoBERTa Base & 0.350 & 0.391 & 0.311 & 0.117 & 0.107 & 0.123 & 0.874 & 0.838 & 0.825 \\
fastText LID-176 & 0.734 & 0.772 & 0.762 & 0.938 & 0.972 & 0.982 & 0.940 & 0.949 & 0.979 \\
ConLID           & 0.764 & 0.732 & 0.762 & 0.955 & 0.923 & 0.971 & 0.987 & 0.946 & 0.969 \\
GlotLID v3       & 0.762 & 0.766 & 0.754 & 0.990 & 0.962 & 0.973 & \textbf{0.995} & 0.969 & 0.978 \\
NLLB LID-218     & 0.766 & 0.766 & 0.762 & \textbf{0.997} & 0.959 & \textbf{0.983} & \textbf{0.995} & \textbf{0.970} & \textbf{0.978} \\
OpenLID-v2       & 0.583 & 0.611 & 0.566 & 0.917 & 0.787 & 0.798 & 0.851 & 0.776 & 0.802 \\
\bottomrule
\end{tabular}%
}
\end{table*}

\paragraph{Zero-Shot Performance.}
When evaluated without modification on our 11-language hybrid test suites (Table \ref{tab:open_world_full}, Left Columns), all six foundation models fail to identify Sinhala-script Pali and Sanskrit, producing \textbf{0.0000 F1}. Because their pretraining label sets map the Sinhala Unicode block exclusively to modern Sinhala, both \texttt{Pali-Sinh} and \texttt{San-Sinh} are classified as Sinhala. This misclassification also lowers modern Sinhala's F1 score to approximately $0.791$. Meanwhile, Sanskrit written in Devanagari (\texttt{San-Deva}) is recognized reliably by most models ($\ge 0.959$), though it drops in OpenLID-v2 ($0.0350$ F1 on FLORES+).

\paragraph{Catastrophic Forgetting in Target-Only Tuning.}
When models are fine-tuned strictly on the three target languages (Table \ref{tab:open_world_full}, Middle Columns), they learn to distinguish Sinhala, Pali, and Sanskrit with high accuracy ($>0.98$ target F1). However, their performance on other languages varies sharply by architecture:
\begin{itemize}
    \item \textbf{Total Collapse in XLM-R:} In XLM-RoBERTa Base, fine-tuning a 3-way classification head with standard cross-entropy alters the shared self-attention layers throughout the network. Without any training signal to anchor its representations for other languages, the model predicts only the three target classes: every non-target language (English, French, German, Tamil, Hindi, Arabic) drops to \textbf{0.0000 F1}. Its overall Macro-F1 collapses to just \textbf{0.1170} on FLORES+, \textbf{0.1074} on CommonLID, and \textbf{0.1234} on WiLI-2018.
    \item \textbf{Stability in Dedicated Models:} In contrast, GlotLID v3, NLLB LID-218, and ConLID preserved their pre-existing language label heads and embedding spaces during continued supervised training. This structural anchor kept their background representations stable, retaining $>0.92$ Macro-F1.
    \item \textbf{Localized Degradation:} Even with continued training, localized degradation still surfaces: fastText LID-176 experiences a drop on English (F1 falling from $0.7213$ to \textbf{0.7074} on FLORES+), and ConLID drops on Arabic ($0.6181$ F1).
\end{itemize}

\subsection{Mitigating Catastrophic Forgetting with Rehearsal}
To prevent catastrophic forgetting, we use continual rehearsal by adding background language examples sampled from the multilingual \textbf{Aya dataset} \cite{ahmed2024aya}. We optimize a balanced replay objective:
\begin{equation}
\mathcal{L} = \alpha \mathcal{L}_{\text{target}} + (1 - \alpha) \mathcal{L}_{\text{rehearsal}}
\end{equation}
where $\alpha = 0.5$ balances updates between target-domain morphology and background language retention.

As shown in Table \ref{tab:open_world_full} (Right Columns), rehearsal delivers clear empirical improvements:
\begin{enumerate}
    \item \textbf{Restoring XLM-RoBERTa:} Rehearsal reverses XLM-R's catastrophic forgetting, raising Macro-F1 by \textbf{$+75.7$ points} from $0.1170$ to \textbf{$0.8740$} on FLORES+ (English rebounds to $0.9646$, Hindi to $0.9975$, and German to $0.9995$), $0.8381$ on CommonLID, and $0.8247$ on WiLI-2018.
    \item \textbf{High Open-World Accuracy:} GlotLID v3 and NLLB LID-218 achieve the strongest overall performance, reaching \textbf{$0.9954$} and \textbf{$0.9948$} Macro-F1 on FLORES+, with consistent $>0.99$ accuracy across both target script classes and global background languages.
    \item \textbf{Resolving Localized Drops:} In ConLID, Arabic F1 recovers from $0.6181$ to \textbf{$0.9877$}, while fastText English F1 improves from $0.7074$ to \textbf{$0.8178$} on FLORES+ and $0.9565$ on WiLI-2018.
\end{enumerate}

\subsection{Inference Speed Comparison}
While both adapted XLM-R and fastText-based foundation models achieve strong Macro-F1, their operational profiles differ substantially:
\begin{itemize}
    \item \textbf{XLM-RoBERTa Base:} Takes approximately \textbf{$18.52\text{ ms}$ per sentence} on an NVIDIA A100 GPU and requires substantial GPU memory ($>1.2\text{ GB}$).
    \item \textbf{Adapted fastText LID-176 / GlotLID v3:} Takes only \textbf{$0.12\text{ ms}$ per sentence on CPU} ($154\times$ faster), with a compact binary size of $126\text{ MB}$.
\end{itemize}
For large-scale ingest pipelines processing millions of documents daily, subword fastText foundation models provide both higher accuracy ($0.9954$ vs.\ $0.8740$ F1) and significantly lower computational cost.

\section{Discussion and Conclusion}
\label{sec:conclusion}

Our findings offer clear guidance for building practical language processing pipelines:

\paragraph{For Digitizing Monastic Archives \& Closed Historical Collections.}
When working with pre-filtered historical collections---such as digitized palm-leaf manuscripts or temple archives known to reside within the Sinhala script---heavy transformer models are unnecessary. Lightweight subword models (like Multinomial Naive Bayes or fastText trained from scratch) reach $\ge 0.995$ Macro-F1 on full sentences and prove far more dependable on isolated catalog headwords ($0.858$ F1) without needing GPU hardware.

\paragraph{For Web Crawlers \& Open-World Pipelines.}
When processing unfiltered global web text, classical models struggle to scale due to feature vocabulary explosion. Here, practitioners should rely on foundation models like GlotLID or fastText. Crucially, \textbf{models should not be fine-tuned solely on target languages} to prevent catastrophic forgetting. Continual adaptation with a balanced multilingual replay buffer (such as Aya) preserves background language accuracy while resolving target-script ambiguities.

\paragraph{Conclusion.}
Treating scripts as exclusive proxies for individual modern languages has left canonical Pali and classical Sanskrit literature completely unrecognized in modern language technologies (0.0\% zero-shot recall). In this work, we created the first dedicated benchmark for these languages in the Sinhala script across closed-world and open-world settings. Our results show that character $n$-gram models best preserve subword morphological distinctions in short fragments, while continual rehearsal with multilingual data effectively prevents catastrophic forgetting in foundation models. We hope these findings assist in preserving and processing historical and liturgical texts within modern multilingual NLP systems.

\section*{Limitations}
We note three primary limitations:
First, roughly 36.8\% of our Sanskrit target corpus (4,684 sentences from DCS) was transliterated from SLP1 Romanization into Sinhala script using \textit{Aksharamukha}. While phonetically and orthographically verified, synthetic transliteration cannot capture every idiosyncratic scribal ligature found in handwritten palm-leaf manuscripts.
Second, our benchmarks test clean digital text; severely degraded historical palm leaves affected by OCR noise require specialized image pre-processing before classification.
Third, while our 11-language hybrid suite samples major Indo-Aryan, Dravidian, and global families, testing across all 200+ world languages remains future work.

\section*{Ethics Statement}
This research respects the cultural and religious sanctity of Theravada canonical texts (\textit{Tipitaka}) and classical Sanskrit treatises. All data sources used in this work (SiPaKosa, SiDiaC-v2.0, SansinNT, DCS, FLORES+, CommonLID, WiLI-2018, Aya) are public-domain or shared under open academic licenses. Our models and benchmarks are dedicated to cultural heritage preservation and improving representation for low-resource languages in global language technology.

\bibliography{custom,anthology-1,anthology-2}

\appendix

\section{Detailed Hybrid Benchmark Results Across 11 Languages}
\label{sec:appendix_tables}

Tables \ref{tab:flores_full}, \ref{tab:common_full}, and \ref{tab:wili_full} report the comprehensive per-class F1 evaluation across all 11 target and background languages on FLORES+, CommonLID, and WiLI-2018 under Zero-Shot, Target-Only Fine-Tuned, and Continual Rehearsal (Aya) conditions.

\begin{table*}[p]
\centering
\scriptsize
\setlength{\tabcolsep}{2.2pt}
\renewcommand{\arraystretch}{0.90}
\caption{FLORES+ Hybrid Benchmark Performance: Per-class F1 across all 11 target and global background languages.}
\label{tab:flores_full}
\resizebox{\textwidth}{!}{%
\begin{tabular}{l ccc cccccccc c}
\toprule
\textbf{Model} & \textbf{Sinh-Sinh} & \textbf{Pali-Sinh} & \textbf{San-Sinh} & \textbf{San-Deva} & \textbf{Eng-Latn} & \textbf{Tam-Taml} & \textbf{Hin-Deva} & \textbf{Ben-Beng} & \textbf{Ara-Arab} & \textbf{Fre-Latn} & \textbf{Ger-Latn} & \textbf{Macro-F1} \\
\midrule
\multicolumn{13}{l}{\textit{Zero-Shot Foundation Models (Pretrained Baseline)}} \\
XLM-R Base         & -- & -- & -- & -- & 0.9985 & -- & 0.1835 & -- & 0.6671 & 0.9990 & 1.0000 & 0.3498 \\
ConLID             & 0.7916 & -- & -- & 0.9985 & 0.9960 & 0.9995 & 0.9970 & 1.0000 & 0.6255 & 1.0000 & 0.9970 & 0.7641 \\
fastText LID-176   & 0.7912 & -- & -- & 0.9594 & 0.7213 & 1.0000 & 0.9629 & 1.0000 & 0.6667 & 0.9873 & 0.9897 & 0.7344 \\
GlotLID v3         & 0.7914 & -- & -- & 0.9950 & 1.0000 & 1.0000 & 0.9926 & 0.9995 & 0.6054 & 1.0000 & 1.0000 & 0.7622 \\
NLLB LID-218       & 0.7912 & -- & -- & 0.9895 & 0.9840 & 1.0000 & 0.9912 & 1.0000 & 0.6664 & 0.9995 & 1.0000 & 0.7656 \\
OpenLID-v2         & 0.7912 & -- & -- & 0.0350 & 0.9995 & 1.0000 & 0.0361 & 0.9970 & 0.5579 & 0.9975 & 1.0000 & 0.5831 \\
\midrule
\multicolumn{13}{l}{\textit{Target-Only Fine-Tuned (3 Target Languages: Sinhala, Pali, Sanskrit)}} \\
XLM-R Base         & 0.1875 & 0.8580 & 0.9989 & -- & 0.0000 & 0.0000 & 0.0000 & 0.0000 & 0.0000 & 0.0000 & 0.0000 & 0.1170 \\
ConLID             & 0.9656 & 0.9754 & 0.9519 & 0.9985 & 0.9960 & 0.9995 & 0.9970 & 1.0000 & 0.6181 & 1.0000 & 0.9970 & 0.9545 \\
fastText LID-176   & 0.9917 & 0.9927 & 0.9839 & 0.9829 & 0.7074 & 1.0000 & 0.9946 & 1.0000 & 0.6667 & 1.0000 & 0.9966 & 0.9379 \\
GlotLID v3         & 0.9956 & 0.9957 & 0.9943 & 0.9955 & 1.0000 & 1.0000 & 0.9931 & 0.9995 & 0.9122 & 1.0000 & 1.0000 & 0.9896 \\
NLLB LID-218       & 0.9969 & 0.9965 & 0.9951 & 0.9905 & 1.0000 & 1.0000 & 0.9926 & 1.0000 & 0.9951 & 0.9995 & 0.9995 & \textbf{0.9969} \\
OpenLID-v2         & 0.8172 & 0.8773 & 0.5252 & 0.0000 & 0.9995 & 1.0000 & 0.9960 & 1.0000 & 0.6667 & 0.9975 & 0.9995 & 0.9172 \\
\midrule
\multicolumn{13}{l}{\textit{Continual Adaptation with Multilingual Rehearsal (Target + Global Languages)}} \\
XLM-R Base         & 0.7310 & 0.9836 & 0.9992 & -- & 0.9646 & 0.9773 & 0.9975 & 0.9208 & 0.5239 & 0.8743 & 0.9995 & 0.8740 \\
ConLID             & 0.9616 & 0.9726 & 0.9409 & 0.9985 & 0.9990 & 0.9995 & 0.9970 & 1.0000 & 0.9877 & 1.0000 & 0.9975 & 0.9868 \\
fastText LID-176   & 0.9766 & 0.9881 & 0.9817 & 0.9885 & 0.8178 & 0.9267 & 1.0000 & 1.0000 & 0.6667 & 0.9956 & 0.9931 & 0.9395 \\
GlotLID v3         & 0.9826 & 0.9939 & 0.9936 & 1.0000 & 0.9907 & 1.0000 & 0.9980 & 0.9995 & 0.9911 & 0.9995 & 1.0000 & \textbf{0.9954} \\
NLLB LID-218       & 0.9852 & 0.9964 & 0.9947 & 1.0000 & 0.9717 & 1.0000 & 0.9995 & 1.0000 & 0.9951 & 1.0000 & 1.0000 & 0.9948 \\
OpenLID-v2         & 0.8674 & 0.8961 & 0.5862 & 0.0000 & 0.8271 & 1.0000 & 1.0000 & 1.0000 & 0.6667 & 0.9208 & 0.9995 & 0.8508 \\
\bottomrule
\end{tabular}%
}
\end{table*}

\begin{table*}[p]
\centering
\scriptsize
\setlength{\tabcolsep}{2.2pt}
\renewcommand{\arraystretch}{0.90}
\caption{CommonLID Hybrid Benchmark Performance: Per-class F1 across all 11 target and global background languages.}
\label{tab:common_full}
\resizebox{\textwidth}{!}{%
\begin{tabular}{l ccc cccccccc c}
\toprule
\textbf{Model} & \textbf{Sinh-Sinh} & \textbf{Pali-Sinh} & \textbf{San-Sinh} & \textbf{San-Deva} & \textbf{Eng-Latn} & \textbf{Tam-Taml} & \textbf{Hin-Deva} & \textbf{Ben-Beng} & \textbf{Ara-Arab} & \textbf{Fre-Latn} & \textbf{Ger-Latn} & \textbf{Macro-F1} \\
\midrule
\multicolumn{13}{l}{\textit{Zero-Shot Foundation Models (Pretrained Baseline)}} \\
XLM-R Base         & -- & -- & -- & -- & 0.9473 & -- & 0.4403 & -- & 0.9955 & 0.9483 & 0.9677 & 0.3908 \\
ConLID             & 0.7916 & -- & -- & 0.9623 & 0.8502 & 0.9818 & 0.9447 & 0.9620 & 0.7683 & 0.8984 & 0.8955 & 0.7323 \\
fastText LID-176   & 0.7910 & -- & -- & 0.8886 & 0.9724 & 0.9643 & 0.9702 & 0.9936 & 0.9947 & 0.9447 & 0.9673 & 0.7715 \\
GlotLID v3         & 0.7914 & -- & -- & 0.9634 & 0.9255 & 0.9878 & 0.9713 & 0.9792 & 0.9631 & 0.9255 & 0.9215 & 0.7662 \\
NLLB LID-218       & 0.7912 & -- & -- & 0.9282 & 0.9240 & 0.9747 & 0.9640 & 0.9863 & 0.9923 & 0.9289 & 0.9410 & 0.7664 \\
OpenLID-v2         & 0.7912 & -- & -- & 0.1169 & 0.8973 & 0.9877 & 0.1749 & 0.9722 & 0.9473 & 0.9052 & 0.9239 & 0.6106 \\
\midrule
\multicolumn{13}{l}{\textit{Target-Only Fine-Tuned (3 Target Languages: Sinhala, Pali, Sanskrit)}} \\
XLM-R Base         & 0.0389 & 0.8141 & 0.9989 & -- & 0.0000 & 0.0000 & 0.0000 & 0.0000 & 0.0000 & 0.0000 & 0.0000 & 0.1074 \\
ConLID             & 0.9656 & 0.9746 & 0.9497 & 0.9624 & 0.8506 & 0.9818 & 0.9447 & 0.9620 & 0.7629 & 0.8986 & 0.8959 & 0.9226 \\
fastText LID-176   & 0.9901 & 0.9924 & 0.9705 & 0.9259 & 0.9335 & 0.9938 & 0.9665 & 0.9933 & 0.9899 & 0.9629 & 0.9748 & 0.9721 \\
GlotLID v3         & 0.9759 & 0.9890 & 0.9921 & 0.9645 & 0.9294 & 0.9818 & 0.9618 & 0.9805 & 0.9610 & 0.9272 & 0.9214 & 0.9622 \\
NLLB LID-218       & 0.9690 & 0.9612 & 0.9929 & 0.9329 & 0.9033 & 0.9875 & 0.9574 & 0.9866 & 0.9936 & 0.9181 & 0.9420 & 0.9586 \\
OpenLID-v2         & 0.7440 & 0.5081 & 0.4720 & 0.0000 & 0.8615 & 0.9756 & 0.9592 & 0.9839 & 0.9912 & 0.8874 & 0.8961 & 0.7865 \\
\midrule
\multicolumn{13}{l}{\textit{Continual Adaptation with Multilingual Rehearsal (Target + Global Languages)}} \\
XLM-R Base         & 0.8243 & 0.9707 & 0.9992 & -- & 0.9606 & 0.9364 & 0.9813 & 0.9611 & 0.9958 & 0.9472 & 0.9684 & 0.8381 \\
ConLID             & 0.9616 & 0.9725 & 0.9386 & 0.9609 & 0.8882 & 0.9818 & 0.9491 & 0.9637 & 0.9884 & 0.9042 & 0.8984 & 0.9461 \\
fastText LID-176   & 0.9766 & 0.9881 & 0.9695 & 0.9303 & 0.9346 & 0.7297 & 0.9752 & 0.9953 & 0.9936 & 0.9638 & 0.9775 & 0.9486 \\
GlotLID v3         & 0.9826 & 0.9939 & 0.9913 & 0.9481 & 0.9616 & 0.9878 & 0.9712 & 0.9822 & 0.9940 & 0.9371 & 0.9132 & 0.9694 \\
NLLB LID-218       & 0.9852 & 0.9911 & 0.9925 & 0.9309 & 0.9497 & 0.9875 & 0.9605 & 0.9855 & 0.9939 & 0.9392 & 0.9509 & \textbf{0.9697} \\
OpenLID-v2         & 0.8133 & 0.7038 & 0.5476 & 0.0000 & 0.9319 & 0.9756 & 0.9794 & 0.9877 & 0.9961 & 0.9183 & 0.9162 & 0.7763 \\
\bottomrule
\end{tabular}%
}
\end{table*}

\begin{table*}[p]
\centering
\scriptsize
\setlength{\tabcolsep}{2.2pt}
\renewcommand{\arraystretch}{0.90}
\caption{WiLI-2018 Hybrid Benchmark Performance: Per-class F1 across target and background languages.}
\label{tab:wili_full}
\resizebox{\textwidth}{!}{%
\begin{tabular}{l ccc cccccccc c}
\toprule
\textbf{Model} & \textbf{Sinh-Sinh} & \textbf{Pali-Sinh} & \textbf{San-Sinh} & \textbf{San-Deva} & \textbf{Eng-Latn} & \textbf{Tam-Taml} & \textbf{Hin-Deva} & \textbf{Ben-Beng} & \textbf{Ara-Arab} & \textbf{Fre-Latn} & \textbf{Ger-Latn} & \textbf{Macro-F1} \\
\midrule
\multicolumn{13}{l}{\textit{Zero-Shot Foundation Models (Pretrained Baseline)}} \\
XLM-R Base         & -- & -- & -- & -- & 0.9527 & -- & 0.1811 & -- & -- & 0.9900 & 0.9899 & 0.3114 \\
ConLID             & 0.7916 & -- & -- & 0.9984 & 0.9448 & 0.9944 & 0.9904 & 0.9435 & -- & 0.9818 & 0.9707 & 0.7616 \\
fastText LID-176   & 0.7912 & -- & -- & 0.9904 & 0.9385 & 0.9950 & 0.9889 & 0.9529 & -- & 0.9774 & 0.9854 & 0.7620 \\
GlotLID v3         & 0.7914 & -- & -- & 0.9914 & 0.9377 & 0.9950 & 0.9173 & 0.9451 & -- & 0.9880 & 0.9754 & 0.7541 \\
NLLB LID-218       & 0.7912 & -- & -- & 0.9909 & 0.9301 & 0.9940 & 0.9889 & 0.9429 & -- & 0.9900 & 0.9904 & 0.7618 \\
OpenLID-v2         & 0.7912 & -- & -- & 0.0276 & 0.9427 & 0.9940 & 0.0134 & 0.9424 & -- & 0.9687 & 0.9817 & 0.5662 \\
\midrule
\multicolumn{13}{l}{\textit{Target-Only Fine-Tuned (3 Target Languages: Sinhala, Pali, Sanskrit)}} \\
XLM-R Base         & 0.2209 & 0.9929 & 0.9989 & -- & 0.0000 & 0.0000 & 0.0000 & 0.0000 & -- & 0.0000 & 0.0000 & 0.1234 \\
ConLID             & 0.9656 & 0.9754 & 0.9471 & 0.9985 & 0.9443 & 0.9944 & 0.9904 & 0.9435 & -- & 0.9818 & 0.9707 & 0.9712 \\
fastText LID-176   & 0.9917 & 0.9927 & 0.9921 & 0.9919 & 0.9381 & 0.9950 & 0.9899 & 0.9507 & -- & 0.9891 & 0.9904 & 0.9822 \\
GlotLID v3         & 0.9956 & 0.9957 & 0.9895 & 0.9975 & 0.9376 & 0.9949 & 0.9107 & 0.9451 & -- & 0.9869 & 0.9738 & 0.9727 \\
NLLB LID-218       & 0.9969 & 0.9964 & 0.9902 & 0.9969 & 0.9452 & 0.9939 & 0.9894 & 0.9429 & -- & 0.9899 & 0.9899 & \textbf{0.9832} \\
OpenLID-v2         & 0.8172 & 0.8685 & 0.5252 & 0.0000 & 0.9203 & 0.9950 & 0.9843 & 0.9446 & -- & 0.9742 & 0.9617 & 0.7977 \\
\midrule
\multicolumn{13}{l}{\textit{Continual Adaptation with Multilingual Rehearsal (Target + Global Languages)}} \\
XLM-R Base         & 0.8473 & 0.9949 & 0.9992 & -- & 0.9012 & 0.9930 & 0.9899 & 0.9733 & -- & 0.9808 & 0.9698 & 0.8247 \\
ConLID             & 0.9616 & 0.9726 & 0.9361 & 0.9990 & 0.9381 & 0.9944 & 0.9904 & 0.9435 & -- & 0.9844 & 0.9712 & 0.9691 \\
fastText LID-176   & 0.9766 & 0.9881 & 0.9895 & 0.9919 & 0.9565 & 0.9749 & 0.9899 & 0.9523 & -- & 0.9847 & 0.9889 & \textbf{0.9793} \\
GlotLID v3         & 0.9826 & 0.9939 & 0.9887 & 0.9985 & 0.9225 & 0.9949 & 0.9853 & 0.9446 & -- & 0.9880 & 0.9759 & 0.9775 \\
NLLB LID-218       & 0.9852 & 0.9964 & 0.9899 & 0.9980 & 0.9068 & 0.9944 & 0.9889 & 0.9424 & -- & 0.9890 & 0.9904 & 0.9781 \\
OpenLID-v2         & 0.8674 & 0.9151 & 0.5864 & 0.0000 & 0.8695 & 0.9950 & 0.9904 & 0.9446 & -- & 0.9827 & 0.9663 & 0.8024 \\
\bottomrule
\end{tabular}%
}
\end{table*}

\section{Dataset Provenance and Source Breakdown}
\label{sec:appendix_data_provenance}

Table \ref{tab:dataset_provenance} summarizes the provenance, orthographic representation, textual domains, sentence distributions, and repository links for all primary sources assembled into our 74,318-sentence target benchmark.

\begin{table*}[t]
\centering
\small
\setlength{\tabcolsep}{5pt}
\caption{Target Corpus Provenance and Source Breakdown for Sinhala, Pali, and Sanskrit in Sinhala Script ($N=74,318$).}
\label{tab:dataset_provenance}
\resizebox{\textwidth}{!}{%
\begin{tabular}{llrlll}
\toprule
\textbf{Language Class} & \textbf{Primary Source / Repository} & \textbf{Sentences} & \textbf{Orthography} & \textbf{Genre / Domain} & \textbf{Public Access Repository} \\
\midrule
\multirow{2}{*}{Sinhala (\texttt{Sinh-Sinh})} 
& SiPaKosa-Sent & 20,418 & Native Sinhala & Buddhist Commentary \& History & \href{https://huggingface.co/datasets/RaniduG/SiPaKosa-Sent}{\texttt{hf.co/datasets/RaniduG/SiPaKosa-Sent}} \\
& Pali--Sinhala Parallel & 11,845 & Native Sinhala & Canonical Translation Exegesis & \href{https://huggingface.co/datasets/sinhala-nlp/pali-sinhala}{\texttt{hf.co/datasets/sinhala-nlp/pali-sinhala}} \\
\midrule
\multirow{2}{*}{Pali (\texttt{Pali-Sinh})} 
& SiDiaC-v.2.0 & 17,472 & Native Sinhala & Canonical Tipitaka \& Atthakatha & \href{https://github.com/NeviduJ/SiDiaC-v.2.0}{\texttt{github.com/NeviduJ/SiDiaC-v.2.0}} \\
& Pali--Sinhala Parallel & 11,845 & Native Sinhala & Canonical Sutta Verses & \href{https://huggingface.co/datasets/sinhala-nlp/pali-sinhala}{\texttt{hf.co/datasets/sinhala-nlp/pali-sinhala}} \\
\midrule
\multirow{3}{*}{Sanskrit (\texttt{San-Sinh})} 
& SansinNT & 7,953 & Native Sinhala & Sanskrit New Testament (27 Books) & \href{https://ebible.org/details.php?id=sansin}{\texttt{ebible.org/details.php?id=sansin}} \\
& DCS (Kalidasa, Asvaghosa) & 4,684 & Transliterated & Classical Kavya \& Shastra & \href{https://github.com/ambuda-org/dcs}{\texttt{github.com/ambuda-org/dcs}} \\
& SiDiaC (Kavacha Sangrahaya) & 101 & Native Sinhala & Protective Mantras \& Stotras & \href{https://github.com/NeviduJ/SiDiaC-v.2.0}{\texttt{github.com/NeviduJ/SiDiaC-v.2.0}} \\
\midrule
\textbf{Total Target Corpus} & & \textbf{74,318} & & & \\
\bottomrule
\end{tabular}%
}
\end{table*}

\section{Baseline Model Architectures and Hyperparameters}
\label{sec:appendix_architectures}

All seven from-scratch baseline models were trained using standardized configurations:
\begin{itemize}
    \item \textbf{Multinomial Naive Bayes:} Scikit-learn \cite{pedregosa2011scikit}; character $n$-gram range $[1, 5]$; sublinear term frequency scaling; Laplace smoothing $\alpha = 0.1$.
    \item \textbf{Linear SVM:} Scikit-learn \texttt{LinearSVC}; character $n$-grams $[1, 5]$; sublinear TF-IDF; squared hinge loss; $C=1.0$; max iterations 2,000.
    \item \textbf{Logistic Regression:} Scikit-learn \texttt{LogisticRegression}; character $n$-grams $[1, 5]$; multinomial cross-entropy; L-BFGS solver; $C=1.0$; max iterations 1,000.
    \item \textbf{XGBoost:} \texttt{XGBClassifier} \cite{chen2016xgboost}; 100 estimators; learning rate $\eta = 0.1$; max depth 6; subsample ratio 0.8; colsample\_bytree 0.8; multi:softprob objective.
    \item \textbf{Char-CNN:} PyTorch implementation \cite{zhang2015character}; character vocab size 200; embedding dim 64; 3 parallel 1D convolution banks with kernel sizes 3, 5, 7 (128 filters each); ReLU; global max-pooling; dropout $p=0.3$; AdamW (lr $10^{-3}$, weight decay $10^{-4}$); batch size 64; 10 epochs with early stopping.
    \item \textbf{Char-BiGRU:} PyTorch implementation; character embedding dim 64; 1-layer bidirectional GRU with 128 hidden units per direction; dropout $p=0.3$; linear head ($256 \to 3$); AdamW (lr $10^{-3}$); batch size 64; 10 epochs.
    \item \textbf{fastText (From Scratch):} C++ binary / Python API \cite{joulin-etal-2017-bag}; character subwords $[3, 6]$; embedding dim 100; lr $\alpha = 0.5$; 10 epochs; hierarchical softmax loss; bucket size 2,000,000.
\end{itemize}

\end{document}
